% Part: first-order-logic
% Chapter: introduction
% Section: first-order-logic

\documentclass[../../../include/open-logic-section]{subfiles}

\begin{document}

\olfileid{fol}{int}{fol}

\olsection{પ્રથમ-ક્રમ તર્કશાસ્ત્ર}

રૂપલક્ષી તર્કશાસ્ત્રના પ્રારંભિક અભ્યાસથી તમે કદાચ પ્રથમ-ક્રમ
તર્કશાસ્ત્રથી પરિચિત હશો.\footnote{હકીકતમાં, અમે લગભગ એવું માની જ
લઈએ છીએ કે તમે પરિચિત છો!{} જો ન હો, તો \emph{forall x}
\citep{Magnus2021} જેવું વધુ પ્રાથમિક પાઠ્યપુસ્તક ફરી જોઈ શકો.} તમે તેને
``પરિમાણાત્મક તર્કશાસ્ત્ર'' અથવા ``વિધેય તર્કશાસ્ત્ર'' તરીકે ઓળખતા
હશો. પ્રથમ-ક્રમ તર્કશાસ્ત્ર સર્વપ્રથમ તો એક રૂપલક્ષી ભાષા છે. એટલે
તેને ચોક્કસ સંકેતભંડોળ છે અને તેનાં પદો આ સંકેતભંડોળમાંથી બનેલી
સંકેતશ્રેણીઓ છે. પરંતુ દરેક સંકેતશ્રેણી માન્ય નથી. માન્ય પદોના જુદા
પ્રકારો છે: પદો, !!{formula}s અને !!{sentence}s. આપણને મુખ્યત્વે
પ્રથમ-ક્રમ તર્કશાસ્ત્રનાં !!{sentence}sમાં રસ છે: તે આપણને અંગ્રેજી
ભાષાનાં વાક્યોનાં રૂપલક્ષી પ્રતિરૂપ આપે છે અને તેમના વિશે આપણે
તર્કશાસ્ત્રીને સામાન્ય રીતે રસ પડે એવા પ્રશ્નો પૂછી શકીએ છીએ. દાખલા
તરીકે:
\begin{itemize}
    \item શું $!B$, $!A$માંથી તાર્કિક રીતે ફલિત થાય છે?
    \item શું $!A$ તાર્કિક રીતે સત્ય, તાર્કિક રીતે અસત્ય કે
નિવાર્ય છે?
    \item શું $!A$ અને $!B$ સમતુલ્ય છે?
\end{itemize}

આ પ્રશ્નો મુખ્યત્વે પ્રથમ-ક્રમ તર્કશાસ્ત્રનાં !!{sentence}sના
``અર્થ'' વિશેના પ્રશ્નો છે. દાખલા તરીકે, $!B$, $!A$માંથી તાર્કિક રીતે
ફલિત થાય છે કે નહિ એ પ્રશ્નનું તત્ત્વચિંતક આ રીતે વિશ્લેષણ કરશે: શું
એવો કોઈ કિસ્સો છે જેમાં $!A$ સત્ય પણ~$!B$ અસત્ય હોય ($!B$,
$!A$માંથી ફલિત થતું નથી), કે પછી $!A$ને સત્ય બનાવતો દરેક કિસ્સો
~$!B$ને પણ સત્ય બનાવે છે ($!B$, $!A$માંથી ફલિત થાય છે)? પરંતુ
``કિસ્સો'' એટલે શું તે હજુ જણાવ્યું નથી---એ \emph{અર્થવિચાર}નું કામ
છે. પ્રથમ-ક્રમ તર્કશાસ્ત્રનો અર્થવિચાર, તત્ત્વચિંતકના ``કિસ્સા'' વિશેના
સહજ ખ્યાલનું અને---આ મહત્ત્વનું છે---કોઈ કિસ્સામાં
!!a{sentence}~$!A$નું \emph{સત્ય હોવું} એટલે શું તેનું ગાણિતિક રીતે
ચોક્કસ પ્રતિરૂપ આપે છે. આપણે વિકસાવવાના આ ગાણિતિક રીતે ચોક્કસ
પ્રતિરૂપને !!a{structure} કહીએ છીએ. ``એમાં સત્ય''ને ચોક્કસ બનાવતા
સંબંધને \emph{સંતોષ}-સંબંધ કહે છે. તેથી આપણે !!{sentence}s~$!A$ અને
!!{structure}s~$\Struct{M}$ માટે ``$!A$,~$\Struct{M}$માં સંતોષાય છે''
(સંકેતોમાં: $\Sat{M}{!A}$) વ્યાખ્યાયિત કરીશું. એ થઈ જાય પછી ``માંથી
ફલિત થાય છે'' અથવા ``તાર્કિક રીતે સત્ય છે'' જેવા બીજા અર્થવિચારી
પદોની પણ ચોક્કસ વ્યાખ્યાઓ આપી શકીએ. આ વ્યાખ્યાઓ વડે, દાખલા તરીકે,
$\lforall[x][(!A(x) \lif !B(x))],
\lexists[x][!A(x)] \Entails \lexists[x][!B(x)]$ છે કે નહિ તે પણ
ગાણિતિક ચોકસાઈથી નક્કી કરી શકાશે. જવાબ અલબત્ત ``હા'' છે. જો તમને
અંગ્રેજી વાક્યોને પ્રથમ-ક્રમ તર્કશાસ્ત્રમાં સંકેતબદ્ધ કરવાની તાલીમ
મળી હોય, તો તમે આને, માનો કે, ``બધી કીડીઓ કીટકો છે, કીડીઓ અસ્તિત્વમાં
છે, તેથી કીટકો અસ્તિત્વમાં છે''ના સંકેતીકરણ તરીકે ઓળખશો. આ દેખીતી
રીતે પ્રમાણભૂત દલીલ છે; તેથી આપણી રૂપલક્ષી ભાષા માટે ``માંથી ફલિત થાય
છે''નું ગાણિતિક પ્રતિરૂપ પણ એ જ જવાબ આપવું જોઈએ.
\sourcecorrection{OLFINT-001}{મૂળની \texttt{first-order-logic.tex},
પંક્તિઓ~51--52માં સાર્વત્રિક સૂત્રનું અંતકૌંસ પહેલા અસ્તિત્વલક્ષી
સૂત્ર પછી મૂકાયું છે; તેથી બંને આધારવિધાનો અનિચ્છિત રીતે એક જ
પરિમાણકના વ્યાપમાં જાય છે. અહીં દરેક આધારવિધાનની કૌંસરચના અલગ કરીને
અભિપ્રેત ફલિતતા પુનઃસ્થાપિત કરી છે.}

રૂપલક્ષી તર્કશાસ્ત્રના પ્રારંભિક અભ્યાસમાંથી તમને યાદ હોય એવો બીજો
વિષય \emph{!!{derivation}s} છે. તમે રૂપલક્ષી તર્કશાસ્ત્રનો પ્રારંભિક
અભ્યાસક્રમ લીધો હોય, તો તમારા શિક્ષકે તમને આવી !!{derivation}s
શોધવાનો અભ્યાસ કરાવ્યો હશે; કદાચ ઉપરની ફલિતતા માન્ય છે એવું બતાવતી
!!a{derivation} પણ. !!{derivation}s આપવાની ઘણી જુદી રીતો છે: તમે
``સ્વાભાવિક નિગમન'' અથવા ``સત્ય-વૃક્ષો'' નામની રીત કરી હશે, પરંતુ
બીજી ઘણી રીતો પણ છે. !!{derivation} તંત્રોનો હેતુ ઉપરના
તર્કશાસ્ત્રીય પ્રશ્નોના જવાબ આપી શકાય એવા સાધનો આપવાનો છે: દાખલા
તરીકે, જે સ્વાભાવિક નિગમન !!{derivation}માં
$\lforall[x][(!A(x) \lif !B(x))]$ અને $\lexists[x][!A(x)]$ આધારવિધાનો
અને $\lexists[x][!B(x)]$ ફલિતવિધાન (છેલ્લી પંક્તિ) હોય, તે
$\lexists[x][!B(x)]$, $\lforall[x][(!A(x) \lif !B(x))]$ અને
$\lexists[x][!A(x)]$માંથી તાર્કિક રીતે ફલિત થાય છે તેની
\emph{ખાતરી} આપે છે.
\sourcecorrection{OLFINT-002}{મૂળની \texttt{first-order-logic.tex},
પંક્તિઓ~67--69માં સાર્વત્રિક આધારવિધાનનું અંતકૌંસ રહી ગયું છે અને
અસ્તિત્વલક્ષી ફલિતવિધાન પછી વધારાનું અંતકૌંસ છે. અહીં બંને સૂત્રોની
કૌંસરચના પૂર્ણ કરી છે.}

પણ એવું શા માટે? પ્રથમ નજરે !!{derivation} તંત્રોને અર્થવિચાર સાથે
કોઈ સંબંધ નથી: રૂપલક્ષી !!{derivation} આપતાં માત્ર સંકેતોને ચોક્કસ
નિયમોને આધીન રીતે ગોઠવવાના હોય છે; તેમાં ``કિસ્સાઓ'' કે ``એમાં
સત્ય''નો ઉલ્લેખ જ નથી. !!{derivation} તંત્રો અને અર્થવિચાર વચ્ચેનો
સંબંધ અધિસિદ્ધાંતલક્ષી તપાસ વડે સ્થાપવો પડે. દાખલા તરીકે, એવી
ગાણિતિક સાબિતી જોઈએ કે આધારવિધાનો
$\lforall[x][(!A(x) \lif !B(x))]$ અને $\lexists[x][!A(x)]$માંથી
$\lexists[x][!B(x)]$ની રૂપલક્ષી !!{derivation} શક્ય હોય ત્યારે, અને
માત્ર ત્યારે જ, $\lforall[x][(!A(x) \lif !B(x))]$ અને
$\lexists[x][!A(x)]$ મળીને $\lexists[x][!B(x)]$ને ફલિત કરે. જોકે આ
કરી શકાય એ પહેલાં વાક્યરચના અને અર્થવિચારની વ્યાખ્યાઓ બરાબર ઘડવા માટે
ઘણું ઝીણવટભર્યું કામ કરવું પડશે.
\sourcecorrection{OLFINT-003}{મૂળની \texttt{first-order-logic.tex},
પંક્તિઓ~81--83માં ફરી સાર્વત્રિક આધારવિધાનનું અંતકૌંસ રહી ગયું છે અને
અસ્તિત્વલક્ષી ફલિતવિધાન પછી વધારાનું અંતકૌંસ છે. અહીં અધિતાર્કિક
દ્વિશરતી દાવામાં ત્રણેય સૂત્રોની કૌંસરચના સુધારી છે.}

\end{document}
